\documentclass[11pt]{article}
\usepackage[margin=1in]{geometry}
\usepackage{amsmath}
\usepackage{graphicx}
\usepackage{hyperref}

\title{Post-typhoon OWT -- Synthesis of pages 29 to 37}
\author{}
\date{}

\begin{document}
\maketitle

\section{Scope and objective}
This document explains the technical architecture required in pages 29 to 37 of
``PostTyphoon\_OWT\_natural\_frequency\_shift\_prediction''. It covers the experimental
scheme, the physics-informed self-calibrating AI, training, validation, and
attention-based XAI.

\section{Experimental design and data acquisition}
\textbf{Goal.} Produce a dense and reproducible dataset linking typhoon loads to
natural-frequency shifts $f_0^{\mathrm{OWT}}$, along with support-condition measurements.

\textbf{Required architecture.}
\begin{itemize}
  \item Scaling laws (Froude) and linear geometry scaling.
  \item Accurate soil reproduction (grain size, Atterberg limits, saturation).
  \item Consolidation and compaction to match in-situ stiffness and strength.
  \item Cyclic loads from local typhoon profiles (OpenFAST/TurbSim).
  \item Precise $f_0^{\mathrm{OWT}}$ shifts from advanced OMA.
\end{itemize}

\section{Transformer-based self-calibrating AI}
\textbf{Core idea.} A transformer integrates a TIM physics core and a slot-attention
module to auto-calibrate soil-pile stiffness.

\subsection{Slot-attention module (self-calibration)}
\textbf{Role.} Learn the initial stiffness matrix $H_{ij}^0$, and a sequence of plastic
stiffness drops $H_{ij}^n$ with their activations $w_n(t)$.

\textbf{Formulation.}
\begin{equation}
H_{ij}(t) = H_{ij}^0 - \sum_{n=1}^{20} w_n(t)\,H_{ij}^n
\end{equation}

\textbf{Slots.} 21 slots total: 1 slot for $H_{ij}^0$ and 20 slots for
$\{w_n(t), H_{ij}^n\}$ pairs.

\textbf{Internal architecture.}
\begin{itemize}
  \item Cross-attention to link slots with encoded inputs.
  \item Self-attention to enforce slot specialization.
  \item Per-slot MLP to decode matrices and activation scalars.
\end{itemize}

\subsection{Pre-training and losses}
\textbf{Data.} High-fidelity FEM data plus TIM-based physics constraints.

\textbf{Physics losses.}
\begin{itemize}
  \item LoT1: frequency consistency (energy conservation).
  \item LoT2: non-negative dissipation and monotonic drops.
  \item Positivity: $H_{ij}(t) \succeq 0$.
\end{itemize}

\textbf{Total loss.}
\begin{equation}
\mathcal{L}_{\mathrm{total}} = \alpha\mathcal{L}_{\mathrm{data}} + \beta\mathcal{L}_{\mathrm{freq}} +
\gamma(\mathcal{L}_{1\mathrm{mon}} + \mathcal{L}_{2\mathrm{mon}}) + \delta\mathcal{L}_{\mathrm{pos}}
\end{equation}

\section{End-to-end transformer}
\textbf{Flow.} Encoder (fixed inputs + load history) $\rightarrow$ self-calibration
module $\rightarrow$ stiffness encoding $\rightarrow$ decoder cross-attention
$\rightarrow$ $H_{ij}(t)$ and $f_0^{\mathrm{OWT}}(t)$.

\textbf{Key points.}
\begin{itemize}
  \item Masked self-attention to avoid future-data leakage.
  \item Small learning rates for slots to preserve robustness.
  \item Physics losses retained during end-to-end training.
\end{itemize}

\section{Testing, validation, and uncertainty}
\textbf{Slot-attention module.}
\begin{itemize}
  \item 85/15 split on ~300 FEM scenarios.
  \item Metrics: MSE, $R^2$, NRMSE, MAPE, bias.
  \item External validation on centrifuge tests (MAPE $< 10\%$, bias $< 5\%$).
\end{itemize}

\textbf{Full model.}
\begin{itemize}
  \item 85/15 split on analytical + scaled experiments.
  \item Out-of-domain validation: reserved tests + field data.
  \item Strong requirement: MAPE $< 1\%$ for $f_0^{\mathrm{OWT}}$.
\end{itemize}

\textbf{Uncertainty.} Gaussian noise on inputs and Monte Carlo dropout to estimate
aleatoric and epistemic uncertainty (target CoV $< 10\%$).

\section{Attention-based XAI}
\textbf{Method.}
\begin{itemize}
  \item Attention heatmaps for slot-attention and full transformer.
  \item Attention roll-out for end-to-end attribution.
  \item Compare to physical expectations (peak loads, stiffness drops).
\end{itemize}

\section{Work performed}
\begin{itemize}
  \item Extracted and analyzed pages 29 to 37.
  \item Structured the required technical architecture and requirements.
  \item Produced this LaTeX synthesis following the requested organization.
\end{itemize}

\end{document}
